% GENERATED FILE -- DO NOT EDIT BY HAND.
% Regenerate with: make thesis-assets  (scripts/thesis_assets.py)
% Editing this file by hand breaks the guarantee that every number in the
% manuscript is derived from committed evidence (ADR-055).

% Source(s): data/finetune/eval-ft-reader-on-reader-val.json; data/finetune/eval-zeroshot-bf16-val.json; data/finetune/eval-zeroshot-qwen-adr059-val.json

\begin{table}[htbp]
  \centering
  \caption{Why the baseline had to be re-measured in the adapter's own numeric precision. The unmatched comparison understates the regression by a factor of 21.}
  \label{tab:precision-confound}
  \begin{tabular}{lrrrl}
    \toprule
    Comparison & Adapter & Baseline & $\Delta$ & Reads as \\
    \midrule
    unmatched (adapter bf16 vs baseline 4-bit) & 0.8246 & 0.8257 & -0.0011 & ``no measurable harm'' \\
    matched (both bf16) & 0.8246 & 0.8480 & -0.0234 & \textbf{a real regression} \\
    \bottomrule
  \end{tabular}

  \vspace{4pt}
  \begin{minipage}{0.95\linewidth}\footnotesize
  \textbf{Protocol note.} Both rows use the same adapter and the same sealed corpus; only the baseline changes. The deployed baseline runs 4-bit on Apple silicon, the adapter runs bf16 on CUDA, and full precision is worth +0.0223 on this corpus. Because the adapter costs roughly what the precision gain is worth, the two nearly cancel: the unmatched difference is small enough to read as neutral. The regression is only visible once the baseline is re-measured in the adapter's own precision, which is why that re-measurement was made a precondition of the study rather than an afterthought.
  \end{minipage}
\end{table}
